% Generated draft: requires mathematical and source review.
\documentclass[11pt]{article}
\usepackage[margin=25mm]{geometry}
\usepackage{fontspec}
\setmainfont{Latin Modern Roman}
\usepackage{amsmath,amssymb,mathtools,unicode-math}
\setmathfont{Latin Modern Math}
\usepackage{enumitem,needspace,xurl,hyperref}
\setlength{\parindent}{0pt}
\setlength{\parskip}{5pt}
\setlength{\emergencystretch}{3em}
\newcommand{\N}{\mathbb{N}}
\newcommand{\Z}{\mathbb{Z}}
\newcommand{\Q}{\mathbb{Q}}
\newcommand{\R}{\mathbb{R}}
\newcommand{\C}{\mathbb{C}}
\newcommand{\F}{\mathbb{F}}
\newcommand{\A}{\mathbb{A}}
\newcommand{\PP}{\mathbb{P}}
\newcommand{\E}{\mathbb{E}}
\newcommand{\eps}{\varepsilon}
\newcommand{\dd}{\,\mathrm{d}}
\begin{document}
\section*{UnsolvedMath ID: 11000005}
\textbf{Source problem code:} catalogue entry
\subsection{Definitions and mathematical statement}
Let $S_g$ be a closed orientable surface of genus $g$ and let $\Gamma_g=\operatorname{MCG}(S_g)$, in the usual orientation-preserving convention. An action of a discrete group $\Gamma$ on a manifold $Y$ is proper when, for every compact $K\subset Y$, the set $\{\gamma\in\Gamma:\gamma K\cap K\ne\varnothing\}$ is finite. This permits finite stabilizers, which matters because mapping class groups have torsion. No cocompactness is required in the following minima:
\[
 n_{\R}(g)=\min\{n:\Gamma_g\text{ acts properly discontinuously on }\R^n\},
\]
\[
 \operatorname{actdim}(\Gamma_g)=\min\{n:\Gamma_g\text{ acts properly discontinuously on a contractible }n\text{-manifold}\}.
\]
Thus $\operatorname{actdim}(\Gamma_g)\le n_{\R}(g)$. The manifold and action regularity conventions should be checked against the source; the excerpt does not specify them.

\textbf{Problem 3.1.} Determine both minima. The source expects them to agree and proposes
\[
 n_{\R}(g)=\operatorname{actdim}(\Gamma_g)=6g-6\qquad(g\ge2).
\]
It states $n(1)=2$. Its upper-bound motivation is the properly discontinuous action on Teichm\"uller space $\mathcal T_g\cong\R^{6g-6}$. The proposed equality is a question, not a theorem asserted here.

\textbf{Suggested lower-bound method.} The source invokes the obstructor-complex technique of [BKK02]. A finite simplicial complex $L$ is used as an $m$-obstructor when it carries the required homological obstruction to embedding in $\R^m$; the full technical definition is not reproduced in the excerpt. Write $S^k_+=S^k\sqcup\{*\}$ for a pointed sphere in the source's disjoint-point sense. Its examples have
\[
 L=S^{k_1}_+*\cdots*S^{k_p}_+,\qquad
 m=k_1+\cdots+k_p+2p-2.
\]
In particular the join of two pointed $0$-spheres is the utilities graph, a $2$-obstructor. Seek a proper expanding map
\[
 F:L\times[0,\infty)\longrightarrow X,
\]
where $X$ is a proper metric space with a proper cocompact isometric $\Gamma$-action. Expansion means that for every pair of disjoint simplices $\sigma,\tau\subset L$,
\[
 \lim_{t\to\infty}d_X\bigl(F(\sigma\times[t,\infty)),F(\tau\times[t,\infty))\bigr)=\infty.
\]
This is separation of whole image tails, not only of points at equal times. The cited method then gives $\operatorname{actdim}(\Gamma)\ge m+2$. For $\Gamma=\Gamma_g$, the proposed choice is the thick part
\[
 X=\mathcal T_g^{\ge\eps}=\{Z\in\mathcal T_g:\text{every closed geodesic of }Z\text{ has length at least }\eps\},
\]
for a sufficiently small positive $\eps$.

The source also reports action-dimension examples: $n(B_k)=2k-3$ for the $k$-strand braid group with $k\ge2$; $n(\Gamma)=\dim(G/K)$ for a lattice in a connected semisimple Lie group $G$, with $K$ maximal compact, citing [BF02a]; and a garbled free-group expression rendered ``$n(Fg\ 2)=2g$.'' The last is plausibly the direct-product example $n((F_2)^g)=2g$, but its exact notation must be checked. These are attributed contextual reports, not independently verified unrestricted theorems. The source says the pointed-sphere joins cover the obstructor applications in [BKK02] and [BF02b].
\subsection{Short English statement}
Find the least dimension of Euclidean space, and of any contractible manifold, admitting a proper mapping-class-group action. The expected answer is $6g-6$ for $g\ge2$; the source proposes a thick-Teichm\"uller-space obstructor construction for the lower bound.
\subsection{Sources}
\begin{itemize}
\item Local UnsolvedMath registry, permanent ID 11000005, source code catalogue entry; no individual proposer or proposal year supplied.
\item Farb, edited volume \emph{Problems on Mapping Class Groups and Related Topics}, Problem 3.1, PDF page 14; \url{https://www.math.uchicago.edu/~farb/papers/mcgbook.pdf}. The registry records PDF-font-structure extraction, accessed 2026-07-29. Its inventory identifies 329 Problem, Conjecture, and Question headings across contributed chapters, disambiguated by PDF page.
\item Source labels [BKK02], [BF02a], and [BF02b], concerning obstructor/action-dimension methods and examples. The supplied record does not include their bibliography entries; the summary associates relevant work with Bestvina--Feighn but does not resolve all labels.
\end{itemize}
\end{document}
